% Section 4: rough weights (Theorem 1.3).
\section{Rough weights}\label{sec:rough}

This section proves Theorem~\ref{thm:rough}. Throughout, $s>1$ is fixed, $v_k:=k^{-s}$ is the pure power, a
superscript $v$ marks quantities of $\sigma_{p,s}$, and $w=ve^\varphi$ with $\varphi$ taking the values $0,\pm\varepsilon$.

The mechanism is the following. By Lemma~\ref{lem:lr}, $\log(b_N^2/b_{N-1}^2)$ and $a_N$ move, to first order, by
the pairing of $\varphi$ with $\Delta^2m_N^v$ and with $\Delta\cE_N^v$, up to a remainder of order $\varepsilon^2$ that holds
for every base measure. For the pure power these two vectors stay of size $\asymp1$ (after the natural scaling) on a
window of indices of size $N$; this follows from the arcsine limit of Lemma~\ref{lem:arcsine}, because a
degenerate vector would force the spectral weights to concentrate near three points
(Lemmas~\ref{lem:window} and~\ref{lem:akernel}). Choosing $\varphi=\pm\varepsilon$ with the sign of that vector on the
window produces a first-order response of size $\asymp\varepsilon$, which the remainder cannot cancel when
$\varepsilon$ is small. Repeating this on widely separated windows gives infinitely many such responses, and the
windows do not interact: the first-order term is evaluated at the pure power, where the other windows lie either in
the saturated head (Corollary~\ref{cor:head}) or beyond $Mn$, where the arcsine limit leaves little mass.

\subsection{Windows and the non-degeneracy of the second difference}\label{sec:window}
A \emph{window} is a pair $W=(\tau_1,\tau_2)$ with $0<\tau_1<\tau_2$; put $W_N:=\N\cap[\tau_1N,\tau_2N]$. It is
\emph{admissible} if $\tau_2>\usat(p)$. Let $Z:=\{-2,0,2\}$, and put
\begin{equation}\label{eq:window-constants}
\begin{gathered}
  I_W:=(\tau_2^{-p},\tau_1^{-p})\cap(0,4c),\qquad \Theta_W:=\{\theta\in(0,\pi):\ 2c(1+\cos\theta)\in I_W\},\\
  \iota_W:=\frac1\pi\int_{\Theta_W}\dist(2\cos\theta,Z)^2\diff\theta,\qquad \delta_W:=\frac{\iota_W}{32},\qquad
  J_W:=\Bigl\lceil\frac2{\delta_W}\Bigr\rceil,\qquad \gamma_W:=\frac{\delta_W}{J_W^2}.
\end{gathered}
\end{equation}
These depend only on $p,\tau_1,\tau_2$. The window is admissible if and only if $I_W\ne\emptyset$, since
$\tau_2^{-p}<4c$ is $\tau_2>\usat$. Then $\iota_W>0$: the map $\theta\mapsto2c(1+\cos\theta)$ takes $(0,\pi)$ onto
$(0,4c)$, so $\Theta_W$ is a non-empty open set, and the integrand vanishes only at $\theta\in\{0,\pi/2,\pi\}$. By
Lemma~\ref{lem:usat}, every window with $\tau_2\ge1$ is admissible.

Since $\arc_{4c}$ is the image of $\diff\theta/\pi$ under $y=2c(1+\cos\theta)$, and $(y-2c)/c=2\cos\theta$,
\begin{equation}\label{eq:iota-arc}
  \iota_W=\int_{I_W}\dist\bigl((y-2c)/c,Z\bigr)^2\diff\arc_{4c}(y).
\end{equation}

\begin{lemma}[window non-degeneracy]\label{lem:window}
Let $W$ be admissible. There is $N_0$ such that every $N\ge N_0$ has some $m\in[N-J_W+1,N]$ with
\[
  \bigl\|\Delta^2m^v_m|_{W_N}\bigr\|_1\ge\gamma_W .
\]
\end{lemma}

The proof is by contradiction. If the second differences are small on $W_N$ for $J_W$ consecutive indices, then so are
the first differences, so $|q_N|$, $|q_{N-1}|$ and $|q_{N-2}|$ nearly agree on $W_N$. The three-term recurrence then
forces the rescaled atoms in $W_N$, weighted by $m_{N-1}$, to lie near the three points $Z$, which contradicts the
arcsine limit, whose mass on $I_W$ is spread out.

\begin{proof}
Write $S:=W_N$, $\delta:=\delta_W$, $J:=J_W$, $\gamma:=\gamma_W$, $D_m:=m_m-m_{m-1}$, and drop the superscript $v$.
Suppose that $\|\Delta^2m_m|_S\|_1<\gamma$ for every $m\in[N-J+1,N]$.

\emph{Step 1: first differences are small on $S$.} For $m\in[N-J,N]$, $D_N-D_m=\sum_{j=m+1}^N\Delta^2m_j$, so
$\|(D_N-D_m)|_S\|_1\le(N-m)\gamma$. Also $m_N-m_{N-J}=\sum_{m=N-J+1}^ND_m$ has $\|\cdot|_S\|_1\le2$. Hence
$J\|D_N|_S\|_1\le2+\gamma J(J-1)/2$, that is, $\|D_N|_S\|_1\le2/J+(J-1)\gamma/2$, and since
$D_{N-1}=D_N-\Delta^2m_N$, also $\|D_{N-1}|_S\|_1\le2/J+(J+1)\gamma/2$. As $2/J\le\delta$ and $\gamma=\delta/J^2$, both
are at most $2\delta$.

\emph{Step 2: the moduli nearly agree on $S$.} Put $e:=|q_{N-2}|-|q_{N-1}|$ and $e':=|q_N|-|q_{N-1}|$. Since
$(\alpha-\beta)^2\le|\alpha^2-\beta^2|$ for $\alpha,\beta\ge0$, $\sum_{k\in S}e_k^2\le\|D_{N-1}|_S\|_1\le2\delta$ and
$\sum_{k\in S}e_k'^2\le\|D_N|_S\|_1\le2\delta$.

\emph{Step 3: the recurrence on $S$.} Fix $k\in S$, let $\epsilon_k:=\sgn(q_{N-1}(k)q_{N-2}(k))$ (with $\sgn0:=1$),
and $t_k:=x_k-a_{N-1}-b_{N-1}\epsilon_k$. We claim
\[
  |q_{N-1}(k)|\cdot\bigl||t_k|-b_N\bigr|\le b_{N-1}|e_k|+b_N|e'_k| .
\]
If $q_{N-1}(k)=0$ there is nothing to prove. Otherwise $q_{N-2}(k)=\epsilon_k\sgn(q_{N-1}(k))|q_{N-2}(k)|$, and
\eqref{eq:rec-vec} at index $N-1$, multiplied by $\sgn(q_{N-1}(k))$, gives
$b_N|q_N(k)|=\bigl|(x_k-a_{N-1})|q_{N-1}(k)|-b_{N-1}\epsilon_k|q_{N-2}(k)|\bigr|=\bigl|t_k|q_{N-1}(k)|-b_{N-1}\epsilon_ke_k\bigr|$.
Hence $\bigl|b_N|q_N(k)|-|t_k||q_{N-1}(k)|\bigr|\le b_{N-1}|e_k|$, and $b_N|q_N(k)|=b_N|q_{N-1}(k)|+b_Ne'_k$, which
proves the claim. Now $||t_k|-b_N|=\dist(t_k,\{\pm b_N\})\ge\dist(x_k-a_{N-1},\tilde Z)$ with
$\tilde Z:=\{\pm b_{N-1}\pm b_N\}$. Put $\bar b:=\max(b_N,b_{N-1})$, $y_k:=(x_k-a_{N-1})/\bar b$ and $Z_N:=\tilde
Z/\bar b$. Squaring the claim, using $(\alpha+\beta)^2\le2\alpha^2+2\beta^2$ and Step~2, and summing over $S$:
\begin{equation}\label{eq:concentration}
  \sum_{k\in S}m_{N-1}(k)\,\dist(y_k,Z_N)^2\le2\bigl(\|e|_S\|_2^2+\|e'|_S\|_2^2\bigr)\le8\delta=\iota_W/4 .
\end{equation}

\emph{Step 4: the arcsine limit contradicts \eqref{eq:concentration} for large $N$.} Put $n:=N-1$ and $Y_k:=n^px_k$.
Then $y_k=\phi_N(Y_k)$ with $\phi_N(Y):=(Y-n^pa_{N-1})/(n^p\bar b)$. By Theorem~\ref{thm:imp-pure},
$n^pa_{N-1}\to2c$ and $n^p\bar b\to c$, so $\phi_N\to\phi$, $\phi(Y):=(Y-2c)/c$, uniformly on compact sets, and
$Z_N\to Z$ in the Hausdorff distance. Since $\dist(\cdot,Z_N)$ is $1$-Lipschitz and differs from $\dist(\cdot,Z)$ by
at most the Hausdorff distance, $\dist(\phi_N,Z_N)^2\to\dist(\phi,Z)^2$ uniformly on compact sets. By
\eqref{eq:iota-arc} and monotone convergence there is a continuous $h$ with $0\le h\le1$, compact support in the open
interval $I_W$, and $\int h\,\dist(\phi,Z)^2\diff\arc_{4c}\ge\iota_W/2$. Now $k\in S$ if and only if $Y_k$ lies in
$[(n/(\tau_2N))^p,(n/(\tau_1N))^p]$, whose endpoints tend to $\tau_2^{-p}$ and $\tau_1^{-p}$; so for large $N$ the
support of $h$ lies inside this interval, and $\mathbf 1_S(k)\ge h(Y_k)$ for all $k$. Hence the left side of
\eqref{eq:concentration} is at least $\int G_N\diff\nu_{n,0}$ with $G_N:=h\,\dist(\phi_N,Z_N)^2$ and $\nu_{n,0}$ as in
Lemma~\ref{lem:arcsine}. The functions $G_N$ converge uniformly on $[0,\infty)$ to the bounded continuous function
$G:=h\,\dist(\phi,Z)^2$, so $\int G_N\diff\nu_{n,0}\to\int G\diff\arc_{4c}\ge\iota_W/2$ by Lemma~\ref{lem:arcsine}. For
$N$ large this contradicts \eqref{eq:concentration}.
\end{proof}

\subsection{The off-diagonal coefficients}\label{sec:rough-b}
Recall $F_N=\log(b_N^2/b_{N-1}^2)$. For an admissible window $W$ put
\[
  M_W:=\max\Bigl(1,\bigl(4c\sin^2(\pi\gamma_W/32)\bigr)^{-1/p}\Bigr),\qquad
  R_W:=\max\bigl(\tau_2/\tau_1,\ \tau_2M_0,\ M_W/\tau_1\bigr).
\]
Since $T_p$ of \eqref{eq:Tp} is non-increasing and $T_p((4c\sin^2(\pi\gamma_W/32))^{-1/p})=\gamma_W/16$, we have
$T_p(M_W)\le\gamma_W/16$.

\begin{theorem}\label{thm:rough-b}
Let $W$ be admissible, let $N_0$ be as in Lemma~\ref{lem:window}, and let $N_1<N_2<\cdots$ be integers with
$N_1\ge\max(N_0,2J_W+2)$, $\tau_2N_i<\tau_1N_{i+1}$, and either
\textup{(A)} $N_{i+1}/N_i\to\infty$, or \textup{(B)} $N_{i+1}>R_WN_i$ for every $i$. For each $i$ pick
$m_i\in[N_i-J_W+1,N_i]$ as in Lemma~\ref{lem:window}. Let $\varepsilon\in(0,\gamma_W/16]$ and
\[
  \varepsilon_k:=\varepsilon\,\sgn\bigl(\Delta^2m^v_{m_i}(k)\bigr)\quad(k\in W_{N_i},\ i\ge1),\qquad
  \varepsilon_k:=0\quad\text{otherwise},
\]
with $\sgn0:=1$. Then $w_k:=k^{-s}e^{\varepsilon_k}$ satisfies $\liminf_iF_{m_i}(w)\ge\varepsilon\gamma_W/2$ in case
\textup{(A)}, and $\liminf_iF_{m_i}(w)\ge\varepsilon\gamma_W/4$ in case \textup{(B)}.
\end{theorem}

\begin{proof}
The windows $W_{N_i}$ are disjoint, so $\varphi:=(\varepsilon_k)$ takes values in $\{0,\pm\varepsilon\}$. Fix $i$, write
$m:=m_i$, and let $H_i:=\bigcup_{j<i}W_{N_j}\subset[1,\tau_2N_{i-1}]$ and
$T_i:=\bigcup_{j>i}W_{N_j}\subset[\tau_1N_{i+1},\infty)$ be the other windows.

\emph{Linearisation.} Lemma~\ref{lem:lr}(ii), applied along the direct path $ve^{t\varphi}$, gives
$F_m(w)\ge F_m(v)+\langle\varphi,\Delta^2m^v_m\rangle-2\varepsilon^2$.

\emph{Own window and cross terms.} By the choice of the signs,
$\langle\varphi,\Delta^2m^v_m\rangle=\varepsilon\|\Delta^2m^v_m|_{W_{N_i}}\|_1+E_i$, where the first term is at least
$\varepsilon\gamma_W$ by Lemma~\ref{lem:window}, and the cross term $E_i:=\langle\varphi\mathbf 1_{H_i\cup T_i},
\Delta^2m_m^v\rangle$ satisfies
$|E_i|\le\varepsilon\bigl(m^v_m+2m^v_{m-1}+m^v_{m-2}\bigr)(H_i\cup T_i)$.

\emph{Head.} Let $m'\in\{m,m-1,m-2\}$; then $N_i-J_W-1\le m'\le N_i$, so $m'\ge N_i/2$. For $k\in H_i$,
$k/m'\le2\tau_2N_{i-1}/N_i$, which tends to $0$ in case (A) and is less than $2\tau_2/R_W\le2/M_0$ in case (B). So, for
$i$ large, $H_i\subset[1,2m'/M_0]$, and Corollary~\ref{cor:head} (with $c_1=1$, $A=s$, $\|v\|_1=\zeta(s)$) gives
$m^v_{m'}(H_i)\le\zeta(s)m'^{s+1}(m'+1)^2e^{-2m'}\to0$.

\emph{Tail.} In case (A) fix $M\ge1$. For $i$ large, $\tau_1N_{i+1}>MN_i\ge Mm'$, so $T_i\subset(Mm',\infty)$ and,
by Lemma~\ref{lem:arcsine}, $\limsup_im^v_{m'}(T_i)\le T_p(M)$; as $M$ is arbitrary, $m^v_{m'}(T_i)\to0$. In case (B),
$\tau_1N_{i+1}>\tau_1R_WN_i\ge M_Wm'$, so $\limsup_im^v_{m'}(T_i)\le T_p(M_W)\le\gamma_W/16$.

\emph{Conclusion.} Hence $\limsup_i|E_i|$ is $0$ in case (A) and at most $4\varepsilon\gamma_W/16$ in case (B). By
Theorem~\ref{thm:imp-pure}, $b_m(v)/b_{m-1}(v)\to1$, so $F_m(v)\to0$. With $\varepsilon\le\gamma_W/16$ we have
$2\varepsilon^2\le\varepsilon\gamma_W/8$, and therefore $\liminf_iF_{m_i}(w)$ is at least
$\varepsilon\gamma_W(1-\frac18)$ in case (A) and $\varepsilon\gamma_W(1-\frac18-\frac14)$ in case (B).
\end{proof}

\begin{proof}[Proof of Theorem~\ref{thm:rough}\textup{(i)}]
Integers $N_i$ as in Theorem~\ref{thm:rough-b} exist (choose them recursively), the weights constructed there are of
the stated form, and $\limsup_nF_n(w)\ge\liminf_iF_{m_i}(w)\ge\varepsilon\gamma_W/4$.
\end{proof}

\subsection{The diagonal coefficients}\label{sec:rough-a}
For the diagonal the relevant vector is $\Delta\cE_n$ of Lemma~\ref{lem:toda}, and its non-degeneracy has a direct proof,
with an explicit constant and a window of two indices. For $z\in\R$ put
\[
  \Upsilon(z):=\min_{\chi\in[0,\pi]}\bigl(|\sin2\chi-z\cos^2\chi|+|z\sin^2\chi-\sin2\chi|\bigr),
\]
and
\begin{equation}\label{eq:ia}
  \ia:=\frac1\pi\int_0^\pi\min\bigl(1,\Upsilon(2\cos\theta)\bigr)\diff\theta .
\end{equation}
The function $\Upsilon$ is $1$-Lipschitz, since the two terms change by at most $|z-z'|\cos^2\chi$ and
$|z-z'|\sin^2\chi$. It vanishes exactly on $Z=\{-2,0,2\}$: with $u=\cos\chi$, $u'=\sin\chi$ the two terms are
$|u(2u'-zu)|$ and $|u'(zu'-2u)|$; if $u=0$ the second is $|z|$, if $u'=0$ the first is $|z|$, and if $uu'\ne0$ both
vanish only if $2u'=zu$ and $zu'=2u$, whence $z^2=4$; conversely the minimum is $0$ for $z\in Z$, and it is attained.
Hence $0<\ia\le1$, because the integrand in \eqref{eq:ia} is continuous, non-negative, and vanishes only at
$\theta\in\{0,\pi/2,\pi\}$.

\begin{lemma}[non-degeneracy of the diagonal response]\label{lem:akernel}
For the pure power,
\[
  \liminf_{N\to\infty}\Bigl(\bigl\|N^p\Delta\cE^v_N\bigr\|_1+\bigl\|N^p\Delta\cE^v_{N+1}\bigr\|_1\Bigr)\ge2c\,\ia .
\]
Consequently, with $\gamma_{\mathrm a}=c\ia/2$, there is $N_0$ such that every $N\ge N_0$ has some $n\in\{N,N+1\}$ with
$\|n^p\Delta\cE^v_n\|_1\ge\gamma_{\mathrm a}$.
\end{lemma}

\begin{proof}
Fix $N$ and write $u:=q_N$, $u':=q_{N+1}$, $B:=b_{N+1}$ (all for $v$). By \eqref{eq:toda-id} at the indices $N$ and
$N+1$, at every atom
\[
  \Delta\cE_N=2Buu'-(x-a_N)u^2,\qquad \Delta\cE_{N+1}=(x-a_{N+1})u'^2-2Buu' .
\]
Put $y_k:=N^px_k$ and $z_k:=(y_k-2c)/c$. Then $N^p\Delta\cE_N=c(2uu'-zu^2)+r_N$ and
$N^p\Delta\cE_{N+1}=c(zu'^2-2uu')+r'_N$, where, by $2|uu'|\le u^2+u'^2$,
\[
  |r_N|+|r'_N|\le\delta_N(u^2+u'^2),\qquad \delta_N:=2|N^pB-c|+|N^pa_N-2c|+|N^pa_{N+1}-2c| ,
\]
and $\delta_N\to0$ by Theorem~\ref{thm:imp-pure}. The expression $|2uu'-zu^2|+|zu'^2-2uu'|$ is homogeneous of degree
$2$ in $(u,u')$, so it is at least $\Upsilon(z)(u^2+u'^2)$ at every atom. Summing over the atoms, with
$\sum_k(u_k^2+u_k'^2)=2$,
\[
  \|N^p\Delta\cE_N\|_1+\|N^p\Delta\cE_{N+1}\|_1\ge c\bigl\langle\min(1,\Upsilon(z)),m_N+m_{N+1}\bigr\rangle-2\delta_N .
\]
By Lemma~\ref{lem:arcsine} with $i=0$ and $i=1$, the images of $m_N$ and $m_{N+1}$ under $k\mapsto y_k$ converge weakly
to $\arc_{4c}$, and $y\mapsto\min(1,\Upsilon((y-2c)/c))$ is bounded and continuous; substituting $y=2c(1+\cos\theta)$,
the pairing tends to $2\ia$. This proves the first claim. For $N\ge N_0$ the sum is then at least $c\ia$, so one of the
two terms is at least $c\ia/2$; and $\|n^p\Delta\cE_n\|_1\ge\|N^p\Delta\cE_n\|_1$ for $n\ge N$.
\end{proof}

\begin{theorem}\label{thm:rough-a}
Let $\varepsilon\in(0,\ia/16]$. Let $N_0$ be as in Lemma~\ref{lem:akernel}, and choose integers $n_1<n_2<\cdots$ with
$\|n_i^p\Delta\cE^v_{n_i}\|_1\ge\gamma_{\mathrm a}$, $n_1\ge\max(N_0,2)$ and $n_{i+1}\ge(i+1)^2M_0n_i$. Put
$W_i:=\N\cap[n_i/M_0,\,in_i]$,
\[
  \varepsilon_k:=(-1)^i\varepsilon\,\sgn\bigl(\Delta\cE^v_{n_i}(k)\bigr)\quad(k\in W_i),\qquad\varepsilon_k:=0\quad
  \text{otherwise},
\]
and $w_k:=k^{-s}e^{\varepsilon_k}$. Then $(-1)^i\bigl(n_i^pa_{n_i}(w)-2c\bigr)\ge\varepsilon\gamma_{\mathrm a}/2-o(1)$ as
$i\to\infty$.
\end{theorem}

\begin{proof}
The indices $n_i$ exist, because every pair $\{N,N+1\}$ with $N\ge N_0$ contains one. The windows are disjoint: the
right end $in_i$ of $W_i$ is less than the left end $n_{i+1}/M_0\ge(i+1)^2n_i$ of $W_{i+1}$. Fix $i$ and put
$A_i:=\N\setminus W_i=[1,n_i/M_0)\cup(in_i,\infty)$. Every other window lies in $A_i$: for $j<i$,
$W_j\subset[1,(i-1)n_{i-1}]$ and $(i-1)n_{i-1}\le(i-1)n_i/(i^2M_0)<n_i/M_0$; for $j>i$, $W_j\subset[n_{i+1}/M_0,\infty)$
and $n_{i+1}/M_0>in_i$.

\emph{Off-window mass.} By Lemma~\ref{lem:toda} and Theorem~\ref{thm:imp-pure}, for $n$ large and every $A\subset\N$,
\[
  n^p\|\Delta\cE^v_n|_A\|_1\le\tfrac12n^p\max(b_n,b_{n+1})(v)\,\bigl(m^v_{n-1}+2m^v_n+m^v_{n+1}\bigr)(A)
  \le c\bigl(m^v_{n-1}+2m^v_n+m^v_{n+1}\bigr)(A).
\]
Let $\omega_i:=c(m^v_{n_i-1}+2m^v_{n_i}+m^v_{n_i+1})(A_i)$. We claim $\omega_i\to0$. For $n'\in\{n_i-1,n_i,n_i+1\}$,
the head $[1,n_i/M_0)$ lies in $[1,2n'/M_0]$, so Corollary~\ref{cor:head} makes its $m^v_{n'}$-mass exponentially
small. For the tail fix $M\ge1$: for $i\ge M$, $(in_i,\infty)\subset(Mn_i,\infty)$, and by Lemma~\ref{lem:arcsine}
with the shifts $-1,0,1$, $\limsup_im^v_{n'}((Mn_i,\infty))\le T_p(M)$. Hence $\limsup_i\omega_i\le4cT_p(M)$ for every
$M$, and $T_p(M)\to0$.

\emph{Linear response.} Lemma~\ref{lem:lr}(iii), applied along the direct path $ve^{t\varphi}$ with
$\varphi=(\varepsilon_k)$, gives
\[
  n_i^p\bigl|a_{n_i}(w)-a_{n_i}(v)-\langle\varphi,\Delta\cE^v_{n_i}\rangle\bigr|
  \le(e^\varepsilon-1-\varepsilon)\,n_i^p\bigl(b_{n_i}+b_{n_i+1}\bigr)(v)=(e^\varepsilon-1-\varepsilon)(2c+o(1)).
\]
By the choice of signs, and because $\varphi$ vanishes on $A_i$ except on the other windows,
\[
  (-1)^in_i^p\langle\varphi,\Delta\cE^v_{n_i}\rangle\ge\varepsilon\,n_i^p\|\Delta\cE^v_{n_i}|_{W_i}\|_1
  -\varepsilon\,n_i^p\|\Delta\cE^v_{n_i}|_{A_i}\|_1\ge\varepsilon(\gamma_{\mathrm a}-2\omega_i).
\]

\emph{Conclusion.} Since $\varepsilon\le\ia/16\le1/16$, $e^\varepsilon-1-\varepsilon\le\frac12\varepsilon^2e^\varepsilon\le
0.54\,\varepsilon^2$, and so $2c(e^\varepsilon-1-\varepsilon)\le1.08\,c\varepsilon\ia/16<\varepsilon\gamma_{\mathrm a}/2$.
With $n^pa_n(v)\to2c$ (Theorem~\ref{thm:imp-pure}) we obtain
$(-1)^i(n_i^pa_{n_i}(w)-2c)\ge\varepsilon\gamma_{\mathrm a}-\varepsilon\gamma_{\mathrm a}/2-o(1)$.
\end{proof}

\begin{proof}[Proof of Theorem~\ref{thm:rough}\textup{(ii)}]
The weights of Theorem~\ref{thm:rough-a} have the stated form, and the values at even and odd $i$ give the bounds for
the $\limsup$ and the $\liminf$.
\end{proof}

\subsection{The averaged law, and remarks}\label{sec:rough-rest}

\begin{proof}[Proof of Theorem~\ref{thm:rough}\textup{(iii)}]
If $e^{-\varepsilon}k^{-s}\le w_k\le e^{\varepsilon}k^{-s}$, then $\sum_kw_k\le e^\varepsilon\zeta(s)$ and $w_k\ge
e^{-\varepsilon}k^{-s}$, so Theorem~\ref{thm:imp-avg} gives \eqref{eq:paperI-logD} and \eqref{eq:paperI-cesaro}. In
case (i), suppose that $n^pb_n\to L\in[0,\infty]$. Then $\log(j^pb_j)\to\log L\in[-\infty,\infty]$, and the weights
$2(n-j)/n^2$ in \eqref{eq:paperI-cesaro} are non-negative with sum $(n-1)/n\to1$, so the left side of
\eqref{eq:paperI-cesaro} tends to $\log L$; hence $L=c$. But then $\log(b_n^2/b_{n-1}^2)=\log(\beta_n^2/\beta_{n-1}^2)
-2p\log\frac n{n-1}\to0$ with $\beta_n:=n^pb_n/c$, which contradicts~(i). In case (ii), the $\limsup$ and the
$\liminf$ of $n^pa_n$ differ by at least $\varepsilon\gamma_{\mathrm a}$.
\end{proof}

\begin{remark}[size of the constants]\label{rem:constants}
The constants of Theorem~\ref{thm:rough} are far from the size of the effect. For $p=2$ and $W=(0.6,1.6)$ a quadrature
gives $\iota_W\approx0.2345$, hence $J_W=273$, $\gamma_W\approx9.8\cdot10^{-8}$ and $R_W\approx1.1\cdot10^8$; for
$p=1$, $\iota_W\approx0.127$ and $\gamma_W\approx1.6\cdot10^{-8}$. Numerically, for the pure power with $p=s=2$, the
norm $\|\Delta^2m_m|_{W_N}\|_1$ is about $1.06$--$1.09$ for every $m\in[N-7,N]$ and $N\in\{100,200,300\}$, and
separation by a factor $4$ already suffices (Section~\ref{sec:numerics}). The loss comes from the contradiction in
Lemma~\ref{lem:window}, which costs the factor $J_W^2$. For the diagonal, a grid quadrature gives $\ia\approx0.314$,
and the observed first-order response $\varepsilon\|n^p\Delta\cE_n^v|_W\|_1$ is about $0.8\,\varepsilon c(p)$ for $p=s=2$,
about ten times the guaranteed $\varepsilon\gamma_{\mathrm a}/2$. None of these numerical values is used in the proofs,
which use only $\iota_W>0$ and $\ia>0$.
\end{remark}

\begin{remark}[the law to accuracy $\varepsilon^2$]\label{rem:accuracy}
Lemma~\ref{lem:lr}(i) shows what a perturbation of size $\varepsilon$ can do to the off-diagonal law. Let $\varphi$ be
bounded with $\frac12\osc\varphi\le\varepsilon$, $w=ve^\varphi$, and $\beta_n:=n^pb_n/c$. Then
$\log\beta_n^2(w)=\langle\varphi,m^v_n-m^v_{n-1}\rangle+r_n+\log\beta_n^2(v)$ with $|r_n|\le\varepsilon^2$ and
$\log\beta_n^2(v)\to0$, so
\[
  \Bigl|\limsup_n|\log\beta_n^2(w)|-\limsup_n|\langle\varphi,m^v_n-m^v_{n-1}\rangle|\Bigr|\le\varepsilon^2 .
\]
Replacing $\varphi$ by $\tau\varphi$ and letting $\tau\downarrow0$, the defect $\limsup_n|\log\beta_n^2(ve^{\tau\varphi})|$
is $o(\tau)$ if and only if $\langle\varphi,m^v_n-m^v_{n-1}\rangle\to0$. So a rough perturbation breaks the law at
first order exactly when it correlates with $m^v_n-m^v_{n-1}$ along infinitely many $n$, which is why the construction
uses sign patterns tied to the spectral weights. For fixed $\varphi$ this leaves a band of width $\varepsilon^2$ that we
cannot close; in particular we do not know whether $\langle\varphi,m_n^v-m_{n-1}^v\rangle\to0$ implies the law for
$ve^\varphi$.
\end{remark}

\begin{remark}[admissibility]\label{rem:admissible}
The condition $\tau_2>\usat$ is necessary for this mechanism. If $\tau_2\le\usat$ and $m-N$ stays bounded, the image
of $W_N$ under $k\mapsto m^px_k$ lies in $[(m/(\tau_2N))^p,\infty)$, whose left end tends to $\tau_2^{-p}\ge4c$; since
$\arc_{4c}$ has no mass beyond $4c$ and no atom, Lemma~\ref{lem:arcsine} gives $m^v_m(W_N)\to0$, so
$\|\Delta^2m^v_m|_{W_N}\|_1\to0$ and the first-order gain disappears. Windows of indices that are saturated in the
sense of Proposition~\ref{prop:imp-eq} carry no spectral weight.
\end{remark}

\begin{remark}\label{rem:rough-ab}
The construction of Theorem~\ref{thm:rough-a} controls $a_n$ only. We do not know whether $b_n$, or the ratio
$a_n/(2b_n)$, converges for the same weights.
\end{remark}
